\documentclass{article}

\usepackage{amsmath,amssymb}
\usepackage{xurl}
\usepackage[hidelinks]{hyperref}
\setlength{\emergencystretch}{2em}

% Shared commands for notes in docs/paper-gaps/.

\DeclareUnicodeCharacter{2080}{\ifmmode{}_{0}\else\textsubscript{0}\fi}
\DeclareUnicodeCharacter{2081}{\ifmmode{}_{1}\else\textsubscript{1}\fi}
\DeclareUnicodeCharacter{2082}{\ifmmode{}_{2}\else\textsubscript{2}\fi}
\DeclareUnicodeCharacter{2099}{\ifmmode{}_{n}\else\textsubscript{n}\fi}

\newcommand{\Lean}{\textsc{Lean}}
\ExplSyntaxOn
\NewDocumentCommand{\leanid}{m}{
  \ifmmode
    \textsf{\detokenize{#1}}
  \else
    \group_begin:
    \urlstyle{sf}
    \tl_set:Nn \l_tmpa_tl {#1}
    \tl_replace_all:Nnn \l_tmpa_tl {\_} {_}
    \exp_args:NV \path \l_tmpa_tl
    \group_end:
  \fi
}
\ExplSyntaxOff
\providecommand{\tr}{\operatorname{tr}}
\newcommand{\tnleanIssueBase}{https://github.com/LionSR/TNLean/issues/}
\newcommand{\tnleanPrBase}{https://github.com/LionSR/TNLean/pull/}
\newcommand{\tnleanDocBase}{https://sirui-lu.com/TNLean/docs/}
\newcommand{\ghissue}[1]{\href{\tnleanIssueBase#1}{GitHub issue \##1}}
\newcommand{\ghpr}[1]{\href{\tnleanPrBase#1}{PR \##1}}
\newcommand{\doclink}[1]{\href{\tnleanDocBase#1.html}{\path{#1}}}

% Dirac notation and the identity operator, as in blueprint/src/macros/common.tex.
% \providecommand keeps note-local definitions authoritative.
\usepackage{dsfont}
\providecommand{\ket}[1]{|#1\rangle}
\providecommand{\bra}[1]{\langle #1|}
\providecommand{\braket}[2]{\langle #1 | #2 \rangle}
\providecommand{\ketbra}[2]{|#1\rangle\!\langle #2|}
\providecommand{\Id}{\mathds{1}}
\providecommand{\one}{\mathds{1}}
\providecommand{\C}{\mathbb{C}}
\providecommand{\MN}[1]{M_{#1}(\C)}
\providecommand{\MD}{M_D(\C)}

% Verdict marker: \gapnote{<kind>}{<status>}. Kind is the policy
% classification (clarification, local-correction, scope-restriction,
% unfaithful, false-source, open-gap); status is open, wip (elimination
% underway), resolved, or historical. Machine-read by the site index;
% prints a verdict line.
\newcommand{\gapnote}[2]{\par\noindent\textbf{Verdict:} #1 (#2).\par}


% Fallbacks keep this note compilable when a different command.tex is found first.
\providecommand{\leanid}[1]{\textsf{\detokenize{#1}}}
\providecommand{\gapnote}[2]{\par\noindent\textbf{Verdict:} #1 (#2).\par}
\providecommand{\ket}[1]{|#1\rangle}
\providecommand{\bra}[1]{\langle #1|}

\title{Allowed Monomials in the Small-Patch Rewrite}
\author{}
\date{2026-10-08}

\begin{document}
\maketitle
\gapnote{scope-restriction}{open}

\section{The source lemma}

Lemma~6.3 of the polynomial PEPS approximation manuscript of September~24,
2026 (\path{05-frames.tex}, lines~188--219) considers old and new one-sheet
encoded frames that agree outside boundedly many affected holes, and $m$
additional square patches satisfying four conditions: (i) the patch outer
squares and the affected holes avoid the outer footprints of the untouched
holes; (ii) the patch inner squares cover every site whose raw owner changes
and the whole outer footprint of every affected hole; (iii) in each patch
outer square, the sites outside the old inner holes have at most two old
owners, and the sites outside the new inner holes have at most two new owners;
(iv) all involved raw owners and affected tag owners belong to a specified
bounded list of parties. The lemma asserts:
\begin{enumerate}
\item the canonical map $M=K_{\rm new}P_m\cdots P_1K_{\rm old}^{\dagger}$ is a
  contraction with
  $\|M\Omega_{\mathcal F_{\rm old}}-\Omega_{\mathcal F_{\rm new}}\|
   \le(m+r_{\rm old})\varepsilon$;
\item for every fixed $a>0$, $M$ has a contractive approximation $M_a$ with
  $\|M_a-M\|\le L^{-a}$ and a polynomial expansion into monomials allowed by
  Theorem~5.2, using only the specified parties.
\end{enumerate}

\section{What the formal statement covers}

The first assertion is formalized with exactly the source hypotheses, for
frames whose tags are listed in any order; the reference error uses only
condition~(i).\footnote{%
\leanid{TNLean.PEPS.EncodedFrame.SmallPatchRewrite.norm_rewrite_le_one_and_refVec_sub_le},
\leanid{TNLean.PEPS.EncodedFrame.SmallPatchRewrite.norm_rewrite_refVec_sub_le}, and, for
arbitrary tag orderings,
\leanid{TNLean.PEPS.EncodedFrame.SmallPatchRewrite.exists_norm_le_one_and_refVec_sub_le_of_perm}.}
For the second assertion the following steps of the proof
(\path{05-frames.tex}, lines~254--341) are formalized: the exact expansion of
$M$ into complete branches and the count of branches; the site-level wire
classification of every branch, for every choice of cylinder terms, including
unequal radii, holes and patches with empty outer sample (whose identity term
selects the empty square), and overlapping patches and holes; the summation of
branchwise errors; the choice of the truncation rank $k$; and the final
rescaling by $1+\tfrac12L^{-a}$.

The residual network of a branch is also formalized. For a branch $(a,s,b)$,
the raw operator $B_aC_sA_b^\dagger$ is an ordered product of rank-one
operators $\ket{k}\bra{b}_S\otimes\mathbb{1}$: the encodings $\ket0\bra{u'_r}$ of the
affected new holes, the patch terms $\ket{w_i}\bra{w_i}$, and the decodings
$\ket{u_t}\bra0$ of the affected old holes. Every such product is the
contraction of the network whose vertices are the bras and the kets of the
factors, kept distinct, and whose wires join the output of a ket at a site to
the next bra containing that site; open input legs are those of bras at sites
reached by no earlier ket, open output legs those of kets at sites met by no
later bra, the legs of the product zero vectors are fixed to $0$, and sites in
no factor are passed through.\footnote{%
\leanid{TNLean.PEPS.SiteChain.chainOp_apply}.}
No wire joins a vertex to itself and every vertex tensor is normalized, so
Lemma~6.2 applies with one group per vertex. An operator read off from a tensor
on the open-leg groups is that tensor, reshaped as a matrix, tensored with the
identity on the untouched sites, so its operator norm is at most the
Hilbert--Schmidt norm of the tensor, with no factor depending on the site
dimensions.\footnote{%
\leanid{TNLean.PEPS.SiteChain.norm_approxOp_le},
\leanid{TNLean.PEPS.SiteChain.exists_chainOp_truncation}.}
Under the four conditions only the patch bras and kets have open legs, so there
are at most $g=2m$ groups; the open inputs of a patch bra have at most two old
owners and the open outputs of a patch ket at most two new owners, all in the
specified list. Hence every branch is approximated within $2m/\sqrt k$ by at most
$k^{2m}$ product terms with coefficients of modulus at most one, and for every
$\delta>0$ (for instance $\delta=L^{-a}$) the rescaled sum
$M_\delta=(1+\tfrac12\delta)^{-1}\sum_\beta\ket a\bra b\otimes
\bigl(\sum_ic_{\beta,i}T_{\beta,i}\bigr)\otimes\mathbb{1}$ is a contraction with
$\|M_\delta-M\|\le\delta$.\footnote{%
\leanid{TNLean.PEPS.EncodedFrame.SmallPatchRewrite.exists_branch_truncation},
\leanid{TNLean.PEPS.EncodedFrame.SmallPatchRewrite.exists_contractive_approx}.}
Each product term $T_{\beta,i}$ is an explicit tensor product of normalized
vectors on the open output groups of the patch kets, normalized covectors on the
open input groups of the patch bras, product zero vectors on the selected hole
squares, and identities.

Not formalized is the last reading of the proof (lines~306--316): that each
product term, together with the tag basis vectors, is the operator of a monomial
allowed by Theorem~5.2. Allowed monomials are available on party layouts:
compositions of contractions on registers of one party, normalized pair sources
and normalized pair effects on registers of two named parties, and exchanges of
tensor factors, as used for Lemma~5.1 and for the bounded-change clauses of
Lemmas~6.5 and~6.6 (\path{polypeps_ownership_change_monomials.tex}). What is
missing is a layout of the frame's registers by owner, with one register per site
and per tag, and the identification of each product term with an allowed party
chain on it: the open-leg group vectors as normalized pair sources and effects
(or private maps), the zero vectors and tag vectors as private maps, and the
identities as untouched registers. The approximation statement is therefore
formalized with the product terms in explicit tensor form, and the clause ``a
polynomial expansion into allowed monomials'' is stated only through that form.

The formal frames record the pairwise disjointness of the physical outer
footprints of holes rather than of the outer squares in the plane; the source
condition implies the formal one, so this is a generalization. The rewrite is
constructed in canonical coordinates in which the tags of the untouched holes
come first in both frames. Since the encodings of holes with disjoint
footprints commute, moving every tag to the new position of its hole is a
relabelling of the tag basis that carries one encoding to the other; along it
the canonical rewrite becomes a contraction between frames with arbitrary tag
orderings, with the same reference error. This is the canonical identification
of tag orderings of Definition~6.1, and it is formalized.

\section{Elimination plan}

Construct the register layout of an encoded frame by owner, with one register
per site and one per tag, and its identification with the canonical
coordinates in which the rewrite acts. On that layout, write each product term
of the branch truncation as an allowed party chain: first the normalized group
covectors on the open input groups (pair effects on the two old owners, or
private maps for one owner), the zero bras and tag bras as private maps, then
the zero kets and tag kets and the normalized group vectors on the open output
groups (pair sources on the two new owners, or private maps), framed by the
untouched registers. The approximation theorem above then yields the second
assertion of Lemma~6.3 with allowed monomials.

\end{document}
